\documentclass[12pt,a4paper]{report}

% ============================================
% PACKAGES
% ============================================
\usepackage[utf8]{inputenc}
\usepackage[T1]{fontenc}
\usepackage[french]{babel}
\usepackage{geometry}
\usepackage{graphicx}
\usepackage{xcolor}
\usepackage{listings}
\usepackage{tcolorbox}
\usepackage{booktabs}
\usepackage{array}
\usepackage{tikz}
\usetikzlibrary{shapes.geometric}
\usepackage{hyperref}
\usepackage{fancyhdr}
\usepackage{titlesec}
\usepackage{enumitem}
\usepackage{amsmath}
\usepackage{amssymb}
\usepackage{float}

% ============================================
% CONFIGURATION
% ============================================
\geometry{margin=2.5cm}
\hypersetup{
    colorlinks=true,
    linkcolor=blue!70!black,
    urlcolor=blue!70!black,
    citecolor=green!50!black
}

% Couleurs personnalisées
\definecolor{primaryblue}{RGB}{41, 98, 255}
\definecolor{secondarygreen}{RGB}{0, 200, 83}
\definecolor{warningorange}{RGB}{255, 152, 0}
\definecolor{dangerred}{RGB}{244, 67, 54}
\definecolor{codebg}{RGB}{245, 245, 245}
\definecolor{codeframe}{RGB}{200, 200, 200}

% Style des listings
\lstset{
    basicstyle=\ttfamily\small,
    backgroundcolor=\color{codebg},
    frame=single,
    rulecolor=\color{codeframe},
    breaklines=true,
    breakatwhitespace=true,
    showstringspaces=false,
    numbers=left,
    numberstyle=\tiny\color{gray},
    keywordstyle=\color{primaryblue}\bfseries,
    commentstyle=\color{secondarygreen},
    stringstyle=\color{dangerred},
    tabsize=4
}

% En-têtes et pieds de page
\pagestyle{fancy}
\fancyhf{}
\fancyhead[L]{\leftmark}
\fancyhead[R]{Modèle Anti-Triche}
\fancyfoot[C]{\thepage}
\renewcommand{\headrulewidth}{0.4pt}
\renewcommand{\footrulewidth}{0.4pt}

% Style des titres
\titleformat{\chapter}[display]
{\normalfont\huge\bfseries\color{primaryblue}}
{\chaptertitlename\ \thechapter}{20pt}{\Huge}

\titleformat{\section}
{\normalfont\Large\bfseries\color{primaryblue!80!black}}
{\thesection}{1em}{}

\titleformat{\subsection}
{\normalfont\large\bfseries\color{primaryblue!60!black}}
{\thesubsection}{1em}{}

% Box personnalisées
\tcbuselibrary{skins,breakable}
\newtcolorbox{infobox}[1][]{
    colback=primaryblue!5,
    colframe=primaryblue,
    title=#1,
    fonttitle=\bfseries,
    breakable
}
\newtcolorbox{successbox}[1][]{
    colback=secondarygreen!10,
    colframe=secondarygreen!70!black,
    title=#1,
    fonttitle=\bfseries,
    breakable
}
\newtcolorbox{warningbox}[1][]{
    colback=warningorange!10,
    colframe=warningorange!70!black,
    title=#1,
    fonttitle=\bfseries,
    breakable
}

% ============================================
% DOCUMENT
% ============================================
\begin{document}

% ============================================
% PAGE DE TITRE
% ============================================
\begin{titlepage}
    \centering
    \vspace*{2cm}
    
    {\Huge\bfseries\color{primaryblue} Système Anti-Triche\\[0.5cm]
    pour Challenges de Code\par}
    
    \vspace{1cm}
    
    {\Large\color{gray} Détection automatique du code généré par IA,\\
    du plagiat et du copier-coller\par}
    
    \vspace{2cm}
    
    \begin{tikzpicture}
        \node[draw=primaryblue, line width=3pt, rounded corners=10pt, 
              inner sep=20pt, fill=primaryblue!5] {
            \begin{minipage}{0.6\textwidth}
                \centering
                \Large\bfseries
                API REST déployée sur\\
                Hugging Face Spaces\\[0.5cm]
                \normalsize\normalfont
                Apprentissage automatique continu\\
                Le modèle devient plus intelligent avec le temps
            \end{minipage}
        };
    \end{tikzpicture}
    
    \vspace{2cm}
    
    {\large
    \begin{tabular}{rl}
        \textbf{URL de l'API :} & \url{https://negzaoui-anti-cheat-coding.hf.space} \\[0.3cm]
        \textbf{GitHub :} & \url{https://github.com/negzaoui16/anti-cheat-coding} \\[0.3cm]
        \textbf{Technologie :} & Python, FastAPI, scikit-learn \\[0.3cm]
        \textbf{Déploiement :} & Docker, Hugging Face Spaces \\
    \end{tabular}
    \par}
    
    \vfill
    
    {\large \today\par}
\end{titlepage}

% ============================================
% TABLE DES MATIÈRES
% ============================================
\tableofcontents
\newpage

% ============================================
% CHAPITRE 1 : INTRODUCTION
% ============================================
\chapter{Introduction}

\section{Contexte}

Dans le cadre des plateformes de challenges de code (coding challenges), il est essentiel de détecter les comportements frauduleux tels que :

\begin{itemize}[leftmargin=2cm]
    \item[\textbf{1.}] \textbf{Code généré par IA} : ChatGPT, GitHub Copilot, Claude, etc.
    \item[\textbf{2.}] \textbf{Plagiat} : Code copié d'autres participants
    \item[\textbf{3.}] \textbf{Copier-coller} : Code collé depuis internet ou d'autres sources
\end{itemize}

\section{Objectif du Système}

Le système anti-triche a pour objectif de :

\begin{successbox}[Objectifs principaux]
\begin{enumerate}
    \item \textbf{Analyser} le code soumis par les participants
    \item \textbf{Calculer} un score de suspicion (0 à 1)
    \item \textbf{Retourner} un verdict : OK, SUSPECT, ou TRICHE PROBABLE
    \item \textbf{Apprendre} automatiquement et devenir plus intelligent avec le temps
    \item \textbf{Ne pas bloquer} automatiquement (décision humaine possible)
\end{enumerate}
\end{successbox}

\section{Caractéristiques Clés}

\begin{table}[H]
\centering
\renewcommand{\arraystretch}{1.5}
\begin{tabular}{>{\bfseries}l l}
\toprule
\textbf{Caractéristique} & \textbf{Description} \\
\midrule
Apprentissage Automatique & Le modèle apprend tout seul \\
API REST & Accessible depuis n'importe quelle application \\
Déploiement Cloud & Hugging Face Spaces (toujours en ligne) \\
CPU uniquement & Pas besoin de GPU \\
Open Source & Code disponible sur GitHub \\
\bottomrule
\end{tabular}
\caption{Caractéristiques principales du système}
\end{table}

% ============================================
% CHAPITRE 2 : ARCHITECTURE
% ============================================
\chapter{Architecture du Système}

\section{Vue d'Ensemble}

Le système est composé de plusieurs modules interconnectés :

\begin{center}
\begin{tikzpicture}[
    box/.style={draw, rounded corners, minimum width=3cm, minimum height=1cm, align=center},
    arrow/.style={->, thick, >=stealth}
]
    % API
    \node[box, fill=primaryblue!20] (api) at (0, 4) {API REST\\(FastAPI)};
    
    % Modules de détection
    \node[box, fill=secondarygreen!20] (ia) at (-4, 2) {Détection IA\\(model\_ia.py)};
    \node[box, fill=secondarygreen!20] (plagiat) at (0, 2) {Détection Plagiat\\(plagiarism.py)};
    \node[box, fill=secondarygreen!20] (copypaste) at (4, 2) {Copier-Coller\\(copy\_paste.py)};
    
    % Scoring
    \node[box, fill=warningorange!20] (scoring) at (0, 0) {Score Final + Verdict\\(scoring.py)};
    
    % Storage
    \node[box, fill=dangerred!20] (storage) at (0, -2) {Stockage + Apprentissage\\(storage.py)};
    
    % Flèches
    \draw[arrow] (api) -- (ia);
    \draw[arrow] (api) -- (plagiat);
    \draw[arrow] (api) -- (copypaste);
    \draw[arrow] (ia) -- (scoring);
    \draw[arrow] (plagiat) -- (scoring);
    \draw[arrow] (copypaste) -- (scoring);
    \draw[arrow] (scoring) -- (storage);
    \draw[arrow, dashed] (storage) to[bend right=60] (ia);
\end{tikzpicture}
\end{center}

\section{Structure des Fichiers}

\begin{lstlisting}[language=bash, caption=Arborescence du projet]
Modele_Cheat/
|-- config.py                 # Configuration globale
|-- requirements.txt          # Dependances Python
|-- Dockerfile                # Deploiement Docker
|-- src/
|   |-- api.py               # API REST FastAPI
|   |-- model_ia.py          # Modele ML (RandomForest)
|   |-- features.py          # Extraction de features
|   |-- copy_paste.py        # Detection copier-coller
|   |-- plagiarism.py        # Detection plagiat (TF-IDF)
|   |-- scoring.py           # Score final + verdict
|   |-- storage.py           # Stockage predictions
|   |-- training.py          # Re-entrainement
|-- models/
|   |-- model_v1.pkl         # Modele entraine
|-- storage/
|   |-- predictions.jsonl    # Historique des analyses
|   |-- human_decisions.jsonl # Decisions humaines
|-- scripts/
    |-- validation_complete_huggingface.py  # Tests
\end{lstlisting}

\section{Description des Modules}

\subsection{Module de Détection IA (model\_ia.py)}

Ce module utilise un classifieur \textbf{RandomForest} de scikit-learn pour prédire si le code a été généré par une IA.

\begin{infobox}[Features Extraites]
\begin{enumerate}
    \item \texttt{mean\_line\_length} : Longueur moyenne des lignes
    \item \texttt{std\_line\_length} : Écart-type de la longueur
    \item \texttt{indentation\_variance} : Variance de l'indentation
    \item \texttt{comment\_ratio} : Ratio de commentaires
    \item \texttt{max\_nesting\_depth} : Profondeur d'imbrication
    \item \texttt{num\_branches} : Nombre de if/for/while
    \item \texttt{avg\_identifier\_length} : Longueur moyenne des identifiants
    \item \texttt{num\_lines} : Nombre de lignes
    \item \texttt{empty\_line\_ratio} : Ratio de lignes vides
\end{enumerate}
\end{infobox}

\subsection{Module de Détection Plagiat (plagiarism.py)}

Utilise \textbf{TF-IDF} (Term Frequency - Inverse Document Frequency) et la \textbf{similarité cosinus} pour comparer le code avec des solutions existantes.

\begin{equation}
\text{Score Plagiat} = \max_{i} \left( \cos(\vec{v}_{\text{code}}, \vec{v}_{\text{solution}_i}) \right)
\end{equation}

\subsection{Module Copier-Coller (copy\_paste.py)}

Détecte le copier-coller basé sur des règles métier :

\begin{itemize}
    \item Blocs de plus de 10 lignes insérés en moins de 500ms
    \item Temps moyen par ligne inférieur à 50ms
\end{itemize}

\subsection{Module de Scoring (scoring.py)}

Combine les trois scores avec une moyenne pondérée :

\begin{equation}
\text{Score Final} = 0.4 \times \text{score\_ia} + 0.3 \times \text{score\_plagiat} + 0.3 \times \text{score\_copier\_coller}
\end{equation}

\begin{table}[H]
\centering
\renewcommand{\arraystretch}{1.3}
\begin{tabular}{c c}
\toprule
\textbf{Score Final} & \textbf{Verdict} \\
\midrule
$< 0.35$ & \textcolor{secondarygreen}{\textbf{OK}} \\
$0.35 - 0.65$ & \textcolor{warningorange}{\textbf{SUSPECT}} \\
$\geq 0.65$ & \textcolor{dangerred}{\textbf{TRICHE PROBABLE}} \\
\bottomrule
\end{tabular}
\caption{Seuils de verdict}
\end{table}

% ============================================
% CHAPITRE 3 : API REST
% ============================================
\chapter{API REST}

\section{URL de Base}

\begin{center}
\Large\texttt{https://negzaoui-anti-cheat-coding.hf.space}
\end{center}

\section{Endpoints Disponibles}

\begin{table}[H]
\centering
\renewcommand{\arraystretch}{1.4}
\begin{tabular}{l l p{7cm}}
\toprule
\textbf{Méthode} & \textbf{Endpoint} & \textbf{Description} \\
\midrule
GET & \texttt{/health} & Vérifie que l'API est en ligne \\
GET & \texttt{/auto\_learn\_status} & État de l'apprentissage automatique \\
POST & \texttt{/analyze} & Analyse du code soumis \\
POST & \texttt{/analyze\_github} & Analyse d'un repo GitHub \\
POST & \texttt{/human\_decision} & Enregistre une décision humaine \\
POST & \texttt{/sync\_storage} & Synchronise vers Hugging Face Hub \\
POST & \texttt{/retrain} & Force un ré-entraînement \\
\bottomrule
\end{tabular}
\caption{Endpoints de l'API}
\end{table}

\section{Endpoint POST /analyze}

\subsection{Requête}

\begin{lstlisting}[language=bash, caption=Exemple d'appel à /analyze]
curl -X POST "https://negzaoui-anti-cheat-coding.hf.space/analyze" \
  -H "Content-Type: application/json" \
  -d '{"code": "def hello():\n    print(\"world\")"}'
\end{lstlisting}

\subsection{Réponse}

\begin{lstlisting}[caption=Réponse de /analyze]
{
  "score_ia": 0.06,
  "score_plagiat": 0.0,
  "copy_paste_detected": false,
  "score_copy_paste": 0.0,
  "final_score": 0.024,
  "verdict": "OK",
  "submission_id": "sub_1234567890"
}
\end{lstlisting}

\section{Endpoint POST /analyze\_github}

\subsection{Requête}

\begin{lstlisting}[language=bash, caption=Exemple d'appel à /analyze\_github]
curl -X POST "https://negzaoui-anti-cheat-coding.hf.space/analyze_github" \
  -H "Content-Type: application/json" \
  -d '{"github_url": "https://github.com/user/repo"}'
\end{lstlisting}

\subsection{Réponse}

\begin{lstlisting}[caption=Réponse de /analyze\_github]
{
  "github_url": "https://github.com/user/repo",
  "owner": "user",
  "repo": "repo",
  "files_analyzed": 14,
  "avg_score_ia": 0.86,
  "final_verdict": "TRICHE PROBABLE",
  "details": [
    {"file": "main.py", "score_ia": 0.95, "verdict": "SUSPECT"},
    {"file": "utils.py", "score_ia": 0.77, "verdict": "SUSPECT"}
  ]
}
\end{lstlisting}

% ============================================
% CHAPITRE 4 : APPRENTISSAGE AUTOMATIQUE
% ============================================
\chapter{Apprentissage Automatique}

\section{Fonctionnement}

Le modèle apprend \textbf{automatiquement} sans intervention humaine :

\begin{center}
\begin{tikzpicture}[
    box/.style={draw, rounded corners, minimum width=4cm, minimum height=1.2cm, align=center, font=\small},
    arrow/.style={->, thick, >=stealth},
    decision/.style={draw, diamond, aspect=2, minimum width=2cm, align=center, font=\small}
]
    % Étapes
    \node[box, fill=primaryblue!20] (analyze) at (0, 6) {POST /analyze\\Code soumis};
    \node[box, fill=secondarygreen!20] (score) at (0, 4) {Calcul score\_ia};
    \node[decision, fill=warningorange!20] (check) at (0, 2) {Score\\extrême ?};
    \node[box, fill=dangerred!20] (autolabel) at (4, 2) {Auto-label\\(0 ou 1)};
    \node[box, fill=primaryblue!20] (counter) at (4, 0) {Compteur +1};
    \node[decision, fill=warningorange!20] (retrain) at (0, 0) {Compteur\\$\geq$ 20 ?};
    \node[box, fill=secondarygreen!20] (train) at (0, -2) {RETRAIN\\automatique};
    \node[box, fill=primaryblue!20] (better) at (0, -4) {Modèle plus\\intelligent !};
    
    % Flèches
    \draw[arrow] (analyze) -- (score);
    \draw[arrow] (score) -- (check);
    \draw[arrow] (check) -- node[above] {Oui} (autolabel);
    \draw[arrow] (check) -- node[left] {Non} (0, 0.5) -- (retrain);
    \draw[arrow] (autolabel) -- (counter);
    \draw[arrow] (counter) -- (retrain);
    \draw[arrow] (retrain) -- node[left] {Oui} (train);
    \draw[arrow] (train) -- (better);
\end{tikzpicture}
\end{center}

\section{Seuils d'Auto-Labellisation}

\begin{table}[H]
\centering
\renewcommand{\arraystretch}{1.5}
\begin{tabular}{c c c}
\toprule
\textbf{Score IA} & \textbf{Action} & \textbf{Label} \\
\midrule
$< 0.15$ & Auto-label & 0 (Humain) \\
$0.15 - 0.85$ & Pas d'action & - \\
$> 0.85$ & Auto-label & 1 (IA) \\
\bottomrule
\end{tabular}
\caption{Seuils d'auto-labellisation}
\end{table}

\section{Configuration}

\begin{lstlisting}[language=python, caption=Configuration de l'apprentissage automatique (config.py)]
# Activer l'apprentissage automatique
AUTO_LEARN_ENABLED = True

# Seuils pour auto-labellisation
AUTO_LABEL_HUMAN_THRESHOLD = 0.15  # score < 0.15 = humain
AUTO_LABEL_IA_THRESHOLD = 0.85     # score > 0.85 = IA

# Re-entrainement apres N decisions
AUTO_RETRAIN_EVERY_N = 20

# Minimum d'echantillons pour retrain
AUTO_RETRAIN_MIN_SAMPLES = 10
\end{lstlisting}

\section{Intervention Humaine (Optionnelle)}

Si le modèle se trompe, un correcteur peut enregistrer la vraie réponse :

\begin{lstlisting}[language=bash, caption=Correction manuelle via /human\_decision]
curl -X POST "https://negzaoui-anti-cheat-coding.hf.space/human_decision" \
  -H "Content-Type: application/json" \
  -d '{
    "submission_id": "sub_123",
    "human_label": 0,
    "comment": "C est un humain"
  }'
\end{lstlisting}

\begin{itemize}
    \item \texttt{human\_label: 0} = C'est un humain (OK)
    \item \texttt{human\_label: 1} = C'est de la triche (IA)
\end{itemize}

% ============================================
% CHAPITRE 5 : RÉSULTATS DES TESTS
% ============================================
\chapter{Résultats des Tests}

\section{Validation Complète}

Le script \texttt{validation\_complete\_huggingface.py} teste tous les aspects du système.

\begin{successbox}[Resultat : 9/9 tests OK]
\texttt{TOTAL: 9 OK / 0 FAIL}

\vspace{0.5cm}
\begin{center}
\fbox{\parbox{0.8\textwidth}{\centering\Large\bfseries VALIDATION REUSSIE\\Le modele Hugging Face fonctionne !}}
\end{center}
\end{successbox}

\section{Tests Effectués}

\begin{table}[H]
\centering
\renewcommand{\arraystretch}{1.4}
\begin{tabular}{c l c}
\toprule
\textbf{\#} & \textbf{Test} & \textbf{Statut} \\
\midrule
1 & GET /health & \textcolor{secondarygreen}{\checkmark OK} \\
2 & GET /auto\_learn\_status & \textcolor{secondarygreen}{\checkmark OK} \\
3 & POST /analyze (code humain) & \textcolor{secondarygreen}{\checkmark OK} \\
4 & POST /analyze (code IA) & \textcolor{secondarygreen}{\checkmark OK} \\
5 & POST /analyze (code vide) & \textcolor{secondarygreen}{\checkmark OK} \\
6 & POST /human\_decision & \textcolor{secondarygreen}{\checkmark OK} \\
7 & POST /sync\_storage & \textcolor{secondarygreen}{\checkmark OK} \\
8 & Test de 12 utilisateurs & \textcolor{secondarygreen}{\checkmark OK} \\
9 & POST /analyze\_github & \textcolor{secondarygreen}{\checkmark OK} \\
\bottomrule
\end{tabular}
\caption{Résumé des tests}
\end{table}

\section{Test des 12 Utilisateurs}

\begin{table}[H]
\centering
\renewcommand{\arraystretch}{1.3}
\begin{tabular}{l c c c}
\toprule
\textbf{Utilisateur} & \textbf{Type} & \textbf{Score IA} & \textbf{Détection} \\
\midrule
User 1 - Débutant & Humain & 0.06 & \textcolor{secondarygreen}{\checkmark} \\
User 2 - Humain erreurs & Humain & 0.05 & \textcolor{secondarygreen}{\checkmark} \\
User 3 - Script simple & Humain & 0.06 & \textcolor{secondarygreen}{\checkmark} \\
User 4 - Humain expérimenté & Humain & 0.37 & \textcolor{secondarygreen}{\checkmark} \\
User 5 - Code brouillon & Humain & 0.16 & \textcolor{secondarygreen}{\checkmark} \\
User 6 - Boucle simple & Humain & 0.02 & \textcolor{secondarygreen}{\checkmark} \\
\midrule
User 7 - ChatGPT style & IA & 0.93 & \textcolor{secondarygreen}{\checkmark} \\
User 8 - Copilot style & IA & 0.67 & \textcolor{secondarygreen}{\checkmark} \\
User 9 - Merge sort IA & IA & 0.93 & \textcolor{secondarygreen}{\checkmark} \\
User 10 - Fibonacci IA & IA & 0.67 & \textcolor{secondarygreen}{\checkmark} \\
User 11 - Palindrome IA & IA & 0.92 & \textcolor{secondarygreen}{\checkmark} \\
User 12 - Classe IA & IA & 0.93 & \textcolor{secondarygreen}{\checkmark} \\
\midrule
\multicolumn{3}{r}{\textbf{Total :}} & \textbf{12/12 (100\%)} \\
\bottomrule
\end{tabular}
\caption{Résultats du test des 12 utilisateurs}
\end{table}

\section{Test GitHub}

\begin{table}[H]
\centering
\renewcommand{\arraystretch}{1.3}
\begin{tabular}{l c c}
\toprule
\textbf{Fichier} & \textbf{Score IA} & \textbf{Verdict} \\
\midrule
config.py & 0.95 & SUSPECT \\
src/api.py & 0.94 & SUSPECT \\
src/copy\_paste.py & 0.99 & SUSPECT \\
src/model\_ia.py & 0.98 & SUSPECT \\
src/features.py & 0.93 & SUSPECT \\
src/plagiarism.py & 0.94 & SUSPECT \\
src/scoring.py & 0.94 & SUSPECT \\
src/storage.py & 0.93 & SUSPECT \\
src/training.py & 0.93 & SUSPECT \\
\midrule
\multicolumn{2}{r}{\textbf{Score moyen :}} & \textbf{0.86} \\
\multicolumn{2}{r}{\textbf{Verdict final :}} & \textcolor{dangerred}{\textbf{TRICHE PROBABLE}} \\
\bottomrule
\end{tabular}
\caption{Analyse du repo GitHub negzaoui16/anti-cheat-coding}
\end{table}

% ============================================
% CHAPITRE 6 : CODE SOURCE
% ============================================
\chapter{Code Source Principal}

\section{API REST (api.py)}

\begin{lstlisting}[language=python, caption=Extrait de src/api.py]
@app.post("/analyze", response_model=AnalyzeResponse)
def analyze_code(request: AnalyzeRequest, background_tasks: BackgroundTasks):
    """Analyse le code et retourne les scores + verdict."""
    code = request.code or ""
    
    # 1. Score IA (modele ML)
    score_ia = predict_ia_score(code)
    
    # 2. Score plagiat (TF-IDF + cosine)
    score_plagiat = plagiarism_checker.score(code)
    
    # 3. Copier-coller (regles)
    cp_detected, score_copy_paste = copy_paste_detected(code)
    
    # 4. Score final et verdict
    result = compute_analysis_result(
        score_ia=score_ia,
        score_plagiat=score_plagiat,
        copy_paste_detected=cp_detected,
        score_copy_paste=score_copy_paste,
    )
    
    # 5. Apprentissage automatique (arriere-plan)
    background_tasks.add_task(_auto_label_and_maybe_retrain, sid, score_ia)
    
    return AnalyzeResponse(**result)
\end{lstlisting}

\section{Extraction de Features (features.py)}

\begin{lstlisting}[language=python, caption=Extrait de src/features.py]
def extract_all_features(code: str) -> dict:
    """Extrait toutes les features du code."""
    lines = code.strip().split("\n")
    mean_len, std_len = extract_line_length_stats(lines)
    indent_var = extract_indentation_variance(lines)
    comment_ratio = extract_comment_ratio(lines)
    max_depth, num_branches, avg_name_len = extract_complexity_indicators(code)
    
    return {
        "mean_line_length": mean_len,
        "std_line_length": std_len,
        "indentation_variance": indent_var,
        "comment_ratio": comment_ratio,
        "max_nesting_depth": max_depth,
        "num_branches": num_branches,
        "avg_identifier_length": avg_name_len,
        "num_lines": len(lines),
        "empty_line_ratio": sum(1 for l in lines if not l.strip()) / len(lines),
    }
\end{lstlisting}

\section{Scoring (scoring.py)}

\begin{lstlisting}[language=python, caption=Extrait de src/scoring.py]
def compute_final_score(score_ia, score_plagiat, score_copy_paste):
    """Score final pondere (entre 0 et 1)."""
    return (
        0.4 * score_ia +
        0.3 * score_plagiat +
        0.3 * score_copy_paste
    )

def get_verdict(final_score):
    """Verdict textuel selon les seuils."""
    if final_score < 0.35:
        return "OK"
    if final_score < 0.65:
        return "SUSPECT"
    return "TRICHE PROBABLE"
\end{lstlisting}

% ============================================
% CHAPITRE 7 : CONCLUSION
% ============================================
\chapter{Conclusion}

\section{Récapitulatif}

\begin{successbox}[Le système anti-triche est opérationnel]
\begin{itemize}
    \item \textbf{9/9 tests} passent avec succès
    \item \textbf{100\%} de détection correcte sur 12 utilisateurs
    \item \textbf{Apprentissage automatique} actif (7 auto-labels créés)
    \item \textbf{API déployée} sur Hugging Face Spaces
    \item \textbf{Analyse GitHub} fonctionnelle
\end{itemize}
\end{successbox}

\section{Points Forts}

\begin{table}[H]
\centering
\renewcommand{\arraystretch}{1.5}
\begin{tabular}{>{\bfseries}l p{8cm}}
\toprule
\textbf{Point Fort} & \textbf{Description} \\
\midrule
Pas de blocage automatique & Le système donne un score, c'est l'humain qui décide \\
Apprentissage continu & Le modèle s'améliore avec le temps sans intervention \\
Léger & Exécutable sur CPU (pas de GPU requis) \\
API publique & Accessible depuis n'importe quelle application \\
Open Source & Code disponible sur GitHub \\
\bottomrule
\end{tabular}
\caption{Points forts du système}
\end{table}

\section{Utilisation Recommandée}

\begin{enumerate}
    \item \textbf{Intégrer l'API} dans votre plateforme de coding challenges
    \item \textbf{Appeler POST /analyze} pour les soumissions de code
    \item \textbf{Appeler POST /analyze\_github} pour les projets GitHub
    \item \textbf{Laisser le modèle apprendre} automatiquement
    \item \textbf{Optionnel :} Corriger les erreurs via POST /human\_decision
\end{enumerate}

\vspace{1cm}

\begin{center}
\Large\bfseries\color{primaryblue}
Le modèle est prêt pour la production !
\end{center}

% ============================================
% FIN DU DOCUMENT
% ============================================
\end{document}
